\documentclass[12pt, a4paper]{article}

% ---------- Encoding & Language ----------
\usepackage[utf8]{inputenc}
\usepackage[T1]{fontenc}
\usepackage[ngerman]{babel}

% ---------- Fonts ----------
\usepackage{lmodern}

% ---------- Layout ----------
\usepackage[
  top=2.5cm,
  bottom=2.5cm,
  left=3cm,
  right=2.5cm
]{geometry}
\setlength{\parskip}{6pt}
\setlength{\parindent}{0pt}

% ---------- Header / Footer ----------
\usepackage{fancyhdr}
\setlength{\headheight}{15pt}
\addtolength{\topmargin}{-3pt}
\pagestyle{fancy}
\fancyhf{}
\fancyhead[L]{\small\leftmark}
\fancyhead[R]{\small\thepage}
\renewcommand{\headrulewidth}{0.4pt}

% ---------- Sections styling ----------
\usepackage{titlesec}

% Section: groß, fett, mit Linie darunter
\titleformat{\section}
  {\large\bfseries}
  {\thesection}{1em}{}
  [\vspace{2pt}\hrule height 0.8pt\vspace{4pt}]
\titlespacing{\section}{0pt}{14pt}{6pt}

% Subsection: normal fett, deutlich sichtbar
\titleformat{\subsection}
  {\normalsize\bfseries}
  {\thesubsection}{1em}{}
  [\vspace{1pt}\hrule height 0.3pt\vspace{2pt}]
\titlespacing{\subsection}{0pt}{10pt}{4pt}

% Subsubsection: normalsize, fett
\titleformat{\subsubsection}
  {\normalsize\bfseries}
  {\thesubsubsection}{1em}{}
\titlespacing{\subsubsection}{0pt}{8pt}{2pt}

% ---------- Boxes ----------
\usepackage[most]{tcolorbox}

\newtcolorbox{definitionbox}[1][]{
  colback=white, colframe=black, coltext=black,
  fonttitle=\bfseries, title={Definition}, #1
}
\newtcolorbox{notebox}[1][]{
  colback=white, colframe=black, coltext=black,
  fonttitle=\bfseries, title={Hinweis}, #1
}
\newtcolorbox{examplebox}[1][]{
  colback=white, colframe=black, coltext=black,
  fonttitle=\bfseries, title={Beispiel}, #1
}

% ---------- TikZ ----------
\usepackage{tikz}
\usetikzlibrary{arrows.meta, positioning, shapes.geometric, fit, backgrounds, decorations.pathreplacing}

% ---------- Math ----------
\usepackage{amsmath, amssymb}

% ---------- Lists ----------
\usepackage{enumitem}
\setlist[itemize]{leftmargin=1.5em, itemsep=2pt}
\setlist[enumerate]{leftmargin=1.5em, itemsep=2pt}

% ---------- Figures ----------
\usepackage{float}

% ---------- Hyperlinks ----------
\usepackage[hidelinks]{hyperref}

% ---------- Misc ----------
\usepackage{microtype}

% =========================================
%                DOCUMENT
% =========================================
\begin{document}

\begin{center}
  {\LARGE\bfseries SWK}\\[6pt]
  {\large Vorlesungsnotizen}\\[4pt]
  {\small Stand: \today}
\end{center}

\vspace{1em}
\tableofcontents
\newpage

% =========================================

% =========================================
\section{Software Konstruktion}
% =========================================

Softwarekonstruktion ist ingenieurmäßiges Vorgehen bei der Erstellung von Software:
planbar, überprüfbar und wiederholbar statt handwerklich und personenabhängig.

\begin{definitionbox}
\textbf{Ingenieurmäßiges Vorgehen} heißt: Anforderungen werden erhoben und
spezifiziert, eine Lösung wird entworfen und dokumentiert, die Umsetzung wird
gemessen und geprüft, und jede dieser Tätigkeiten hinterlässt ein nachvollziehbares
Ergebnis.
\end{definitionbox}

\subsection{Strukturelle und funktionale Qualität}

\begin{itemize}
  \item \textbf{Funktionale Qualität:} Tut die Software das Richtige?
        Merkmale: Korrektheit, Zuverlässigkeit, Sicherheit, Performanz.
        Sichtbar für den Nutzer, prüfbar durch Tests gegen Anforderungen.
  \item \textbf{Strukturelle Qualität:} Ist die Software gut gebaut?
        Merkmale: Wartbarkeit, Verständlichkeit, Testbarkeit, Wiederverwendbarkeit.
        Unsichtbar für den Nutzer, prüfbar durch Metriken, Reviews und statische
        Analyse.
\end{itemize}

Der Unterschied ist prüfungsrelevant: funktionale Qualität wird gegen die
Spezifikation getestet, strukturelle Qualität gegen Standards und Metriken gemessen.

\begin{notebox}
Schlechte strukturelle Qualität kostet nicht sofort. Sie kostet bei jeder Änderung,
und der Effekt ist kumulativ. Genau deshalb braucht sie automatisierte Messung: sie
meldet sich nicht von allein.
\end{notebox}

% =========================================
\section{Anforderungen und Spezifikation}
% =========================================

\subsection{Arten von Anforderungen}

\begin{itemize}
  \item \textbf{Funktionale Anforderungen:} Was das System tun soll.
        \glqq{}Das System berechnet den Rabatt aus Warenkorbwert und Kundenstatus.\grqq{}
  \item \textbf{Nichtfunktionale Anforderungen (Qualitätsattribute):} Wie gut es das
        tun soll. Antwortzeit, Verfügbarkeit, Skalierbarkeit, Sicherheit.
  \item \textbf{Randbedingungen (Constraints):} Vorgaben, die nicht verhandelbar sind.
        Zielplattform, gesetzliche Auflagen, vorhandene Systeme.
\end{itemize}

Nichtfunktionale Anforderungen treiben die Architektur; funktionale Anforderungen
treiben die Implementierung. Deshalb beginnt ADD (Abschnitt~\ref{sec:add}) mit den
Qualitätsattributen und nicht mit den Use Cases.

\subsection{Gute Anforderungen}

Eine Anforderung ist brauchbar, wenn sie \textbf{überprüfbar} ist. Die übliche
Checkliste:

\begin{center}
\begin{tabular}{ll}
\hline
\textbf{Kriterium} & \textbf{Gegenbeispiel} \\
\hline
eindeutig      & \glqq{}Das System soll schnell sein.\grqq{} \\
vollständig    & \glqq{}Fehler werden behandelt.\grqq{} (welche? wie?) \\
widerspruchsfrei & zwei Anforderungen fordern Gegenteiliges \\
überprüfbar    & \glqq{}benutzerfreundlich\grqq{} ohne Messvorschrift \\
verfolgbar     & Anforderung ohne Quelle und ohne Test \\
\hline
\end{tabular}
\end{center}

Aus \glqq{}schnell\grqq{} wird überprüfbar: \emph{95\,\% der Suchanfragen werden in
unter 200\,ms beantwortet, gemessen serverseitig bei 100 gleichzeitigen Nutzern.}

\subsection{User Stories und Use Cases}

\textbf{User Story:} \emph{Als \textless{}Rolle\textgreater{} möchte ich
\textless{}Ziel\textgreater{}, um \textless{}Nutzen\textgreater{}.}
Kurz, verhandelbar, mit Akzeptanzkriterien. Sie ist eine Erinnerung an ein Gespräch,
keine Spezifikation.

\textbf{Use Case:} Ablauf mit Akteur, Vorbedingung, Standardablauf, Alternativabläufen
und Nachbedingung. Deutlich formaler und geeignet, wenn das Zusammenspiel vieler
Schritte das Schwierige ist.

\subsection{Logik: Erfüllbarkeit, Widerspruch, Implikation}
\label{sec:logik}

Anforderungen lassen sich als aussagenlogische Formeln über Booleschen Variablen
schreiben. Drei Begriffe, die in der Verifikation wiederkehren:

\begin{definitionbox}
Seien $\varphi, \psi$ Formeln in KNF.
\begin{itemize}
  \item \textbf{Erfüllbarkeit (SAT):} $\varphi$ ist erfüllbar, wenn eine Belegung
        $v$ existiert mit $\beta[\varphi, v] = 1$.
  \item \textbf{Widerspruch:} $\varphi$ und $\psi$ sind widersprüchlich, wenn
        $(\varphi \wedge \psi)$ nicht erfüllbar ist.
  \item \textbf{Implikation:} $\varphi \models \psi$ gilt, wenn
        $(\varphi \wedge \neg\psi)$ nicht erfüllbar ist.
\end{itemize}
\end{definitionbox}

\begin{notebox}
Der letzte Punkt ist der Schlüssel zum gesamten Verifikationsteil der Vorlesung.
\emph{Jede} Implikationsfrage wird beantwortet, indem man ein Modell der Negation
sucht. Findet man eines, ist es ein Gegenbeispiel; findet man keines, gilt die
Implikation.
\end{notebox}

% =========================================
\section{Vorgehensmodelle und Prozessreife}
% =========================================

\subsection{Wasserfall und V-Modell}

Das \textbf{Wasserfallmodell} ordnet die Phasen streng sequenziell: Anforderungen,
Entwurf, Implementierung, Test, Betrieb. Jede Phase endet mit einem Dokument, und
Rücksprünge sind teuer.

Das \textbf{V-Modell} biegt den Wasserfall zu einem V und ordnet jeder
konstruktiven Phase eine Teststufe zu.

\begin{figure}[H]
\centering
\begin{tikzpicture}[
  box/.style={draw, minimum width=3.1cm, minimum height=0.75cm, align=center,
              font=\small},
  >=Stealth, node distance=0.35cm
]
  \node[box] (anf)  at (0, 4.2)  {Anforderungen};
  \node[box] (grob) at (1.5, 2.8){Grobentwurf};
  \node[box] (fein) at (3.0, 1.4){Feinentwurf};
  \node[box] (impl) at (4.5, 0.1){Implementierung};

  \node[box] (unit) at (6.0, 1.4){Modultest};
  \node[box] (integ)at (7.5, 2.8){Integrationstest};
  \node[box] (abn)  at (9.0, 4.2){Abnahmetest};

  \draw[->] (anf) -- (grob);
  \draw[->] (grob) -- (fein);
  \draw[->] (fein) -- (impl);
  \draw[->] (impl) -- (unit);
  \draw[->] (unit) -- (integ);
  \draw[->] (integ) -- (abn);

  \draw[dashed, <->] (anf)  -- node[above, font=\scriptsize] {validiert} (abn);
  \draw[dashed, <->] (grob) -- node[above, font=\scriptsize] {validiert} (integ);
  \draw[dashed, <->] (fein) -- node[above, font=\scriptsize] {validiert} (unit);
\end{tikzpicture}
\caption{V-Modell: jede Entwurfsstufe hat ihre Teststufe}
\end{figure}

Die waagerechten Linien sind der eigentliche Inhalt: der Abnahmetest prüft gegen die
Anforderungen, der Integrationstest gegen den Grobentwurf, der Modultest gegen den
Feinentwurf. Ein Test ohne zugehöriges Dokument prüft gegen nichts.

\subsection{Agiles Vorgehen und Scrum}

Scrum ersetzt die lange Phasenkette durch kurze, vollständige Zyklen.

\begin{center}
\begin{tabular}{lll}
\hline
\textbf{Rollen} & \textbf{Artefakte} & \textbf{Ereignisse} \\
\hline
Product Owner & Product Backlog & Sprint Planning \\
Scrum Master  & Sprint Backlog  & Daily Scrum \\
Entwicklungsteam & Increment    & Sprint Review \\
              & Definition of Done & Sprint Retrospective \\
\hline
\end{tabular}
\end{center}

Der Product Owner verantwortet \emph{was}, das Team verantwortet \emph{wie}, der
Scrum Master verantwortet, dass der Prozess funktioniert. Diese Trennung ist der
häufigste Prüfungspunkt.

\begin{notebox}
Scrum ist kein Vorgehensmodell für Entwurf, sondern für Steuerung. Was gebaut wird,
entscheidet der Backlog; \emph{wie} es strukturiert wird, sagt Scrum nicht. Deshalb
schließen sich Scrum und Architekturarbeit nicht aus.
\end{notebox}

\subsection{CMMI: Prozessreife}

CMMI misst nicht die Software, sondern den \emph{Prozess}, der sie herstellt. Fünf
Reifegrade:

\begin{center}
\begin{tabular}{clp{7.4cm}}
\hline
\textbf{Grad} & \textbf{Name} & \textbf{Kennzeichen} \\
\hline
1 & Initial     & Erfolg hängt an einzelnen Personen, Prozesse ad hoc \\
2 & Managed     & Projekte werden geplant, verfolgt und gesteuert \\
3 & Defined     & Ein organisationsweiter Standardprozess existiert \\
4 & Quantitatively Managed & Der Prozess wird quantitativ gemessen und gesteuert \\
5 & Optimizing  & Der Prozess wird auf Basis der Messung systematisch verbessert \\
\hline
\end{tabular}
\end{center}

Der Sprung von 3 auf 4 ist der interessante: ab dort wird der Prozess mit Zahlen
gesteuert, was ohne Grad 3 nicht geht, weil es vorher keinen gemeinsamen Prozess
gibt, über den man Zahlen bilden könnte.

% =========================================
\section{Schätzen und Messen}
% =========================================

\subsection{Zwei- und Dreipunktschätzung}

Der Aufwand eines Arbeitspakets ist eine Zufallsgröße mit unbekannter Dichte. Man
schätzt drei Werte und \emph{nimmt eine Verteilung an}:

\begin{definitionbox}
Mit kleinstem Wert $a$, mittlerem Wert $c$ und größtem Wert $b$:
\[
E(x) = \frac{a + r\cdot c + b}{2 + r}, \qquad
V(x) = \frac{(b-a)^2}{u^2}, \qquad
S(x) = \frac{b-a}{u}
\]
\end{definitionbox}

\begin{center}
\begin{tabular}{lcc}
\hline
\textbf{Annahme} & $r$ & $u$ \\
\hline
Zweipunktschätzung & 0 & 6 \\
PERT (Beta-Näherung) & 4 & 6 \\
Beta ($p=q=3$, quadratisch) & 4 & 5{,}29 \\
Beta ($p=q=4$, kubisch) & 6 & 6 \\
Beta ($p=4$, $q=2$) & 4 & 5{,}6 \\
\hline
\end{tabular}
\end{center}

Damit ist die Zweipunktschätzung $E = \frac{a+b}{2}$ und die Dreipunktschätzung
$E = \frac{a + 4c + b}{6}$, beide mit $S = \frac{b-a}{6}$.

\subsection{Aggregation über Arbeitspakete}

\[
E = \sum_i E_i, \qquad S = \sqrt{\sum_i S_i^2}
\]

Erwartungswerte addieren sich, Standardabweichungen \emph{nicht}: es addieren sich die
Varianzen. Fünf Pakete mit je $S_i = 6$ ergeben zusammen $S = 13{,}4$ und nicht $30$.

\begin{notebox}
Das gilt nur bei \textbf{unabhängigen} Paketen. Ein gemeinsamer Engpass, ein
gemeinsamer Zulieferer oder eine Person, die in allen Paketen steckt, zerstört die
Annahme, und keine Formel bemerkt das.
\end{notebox}

\subsubsection*{Beispiel: Spieleprojekt (Übungsblatt 2)}

\begin{center}
\begin{tabular}{lrrr|rr}
\hline
\textbf{Komponente} & $a$ & $c$ & $b$ & $E$ (3-Pkt.) & $S$ \\
\hline
Datenverwaltung   &  8 &  20 &  25 &  18{,}83 &  2{,}83 \\
Gamelogik         & 30 &  64 & 128 &  69{,}00 & 16{,}33 \\
Grafikengine      & 42 & 121 & 180 & 117{,}67 & 23{,}00 \\
Gamepad-Interface &  8 &  16 &  20 &  15{,}33 &  2{,}00 \\
Security          & 10 &  30 &  45 &  29{,}17 &  5{,}83 \\
\hline
\textbf{Summe}    &    &     &     & \textbf{250{,}00} & \textbf{29{,}01} \\
\hline
\end{tabular}
\end{center}

Zweipunktschätzung: $E = 248{,}00$ bei identischem $S = 29{,}01$. Das ist kein Zufall:
$S$ hängt nur von der Spannweite ab, nicht vom mittleren Wert.

\subsection{Velocity}

\begin{definitionbox}
\textbf{Velocity} ist die Anzahl Story Points, die ein Team pro Sprint
\emph{fertigstellt}. Story Points messen Umfang, nicht Zeit; Velocity ist die einzige
Brücke zwischen beiden, und sie wird gemessen, nicht geschätzt.
\end{definitionbox}

Regeln aus der Vorlesung:

\begin{itemize}
  \item Nur \textbf{fertige} Stories zählen. Halb fertige Arbeit bringt null Punkte.
  \item Nach etwa drei Sprints pendelt sich die Velocity ein.
  \item Danach plant man mit dem \textbf{Median} der letzten Sprints, nicht mit dem
        Mittelwert: der Median ist robust gegen einen einzelnen schlechten Sprint.
\end{itemize}

\subsubsection*{Beispiel (Übungsblatt 2, Aufgabe 2.2)}

\begin{center}
\begin{tabular}{crrl}
\hline
\textbf{Sprint} & \textbf{geplant} & \textbf{tatsächlich} & \textbf{fertiggestellt} \\
\hline
1 & 15 & 11 & U2, U3, U4 \\
2 & 16 & 11 & U1, U5, U7 \\
3 & 15 & 10 & U6, U8 \\
4 &  9 &  5 & U9 \\
5 &  4 &  4 & U10 \\
\hline
\end{tabular}
\end{center}

Median der letzten drei Sprints: $5$. Insgesamt wurden $41$ von $59$ zugesagten
Punkten geliefert, das Team übernimmt sich also systematisch. Zwei typische Ursachen
für sinkende Velocity: wachsende technische Schulden und Aufgaben, die nicht im
Sprint Backlog stehen (Support, Störungen).

\subsection{Planning Poker}

Alle schätzen gleichzeitig, die Extreme begründen sich, es wird neu geschätzt. Das
Ergebnis ist der Konsens, das Produkt ist die Diskussion.

Die Kartenfolge ist bewusst nicht linear: $1, 2, 3, 5, 8, 13, 21$. Große Stories
lassen sich nicht genauer schätzen als kleine, und die wachsenden Abstände sagen
das ehrlich, statt eine Unterscheidung zwischen 17 und 18 anzubieten.

\textbf{Triangulation:} Schätzungen werden gegeneinander geprüft, nicht gegen Stunden.
Ist die 2-Punkte-Story wirklich doppelt so groß wie die 1-Punkte-Story? Sind beide
5-Punkte-Stories ähnlich groß?

\subsection{Messen im Betrieb: GQM und KPI}

\begin{definitionbox}
\textbf{GQM (Goal -- Question -- Metric)} nach Basili und Rombach: erst das Ziel, dann
die Fragen, die man beantworten muss, und erst dann die Metriken, die diese Fragen
beantworten. Nie in der umgekehrten Richtung.
\end{definitionbox}

\begin{figure}[H]
\centering
\begin{tikzpicture}[
  node/.style={draw, rounded corners, minimum height=0.7cm, align=center,
               font=\small, inner sep=4pt},
  >=Stealth
]
  \node[node] (g) at (0, 2.4) {\textbf{Goal:} Conversion Rate im Onboarding erhöhen};
  \node[node] (q1) at (-4.2, 1.1) {Q1: Wie viele\\ schließen ab?};
  \node[node] (q2) at (0, 1.1)    {Q2: Wie lange\\ dauert es?};
  \node[node] (q3) at (4.2, 1.1)  {Q3: Wo brechen\\ sie ab?};
  \node[node] (m1) at (-4.2, -0.4){M: Abschlussquote};
  \node[node] (m2) at (0, -0.4)   {M: Median-Dauer};
  \node[node] (m3) at (4.2, -0.4) {M: Schritt mit\\ größtem Abbruch};

  \draw[->] (g) -- (q1);  \draw[->] (g) -- (q2);  \draw[->] (g) -- (q3);
  \draw[->] (q1) -- (m1); \draw[->] (q2) -- (m2); \draw[->] (q3) -- (m3);
\end{tikzpicture}
\caption{GQM-Baum: die Pfeilrichtung ist die Methode}
\end{figure}

Ein \textbf{KPI} ist eine Kennzahl, an die eine \emph{Entscheidung} gekoppelt ist.
Eine Kennzahl, bei der niemand anders handeln würde, ist Dekoration.

\begin{notebox}
\textbf{Goodharts Gesetz:} Sobald eine Kennzahl zum Ziel wird, hört sie auf, ein gutes
Maß zu sein. Typisches Muster: die Kennzahl steigt stetig, das eigentliche Ergebnis
bewegt sich nicht.
\end{notebox}

\subsection{A/B-Tests}

Monitoring sagt, \emph{was} passiert ist. Ein Experiment sagt, ob die eigene Änderung
es \emph{verursacht} hat.

Der Stichprobenumfang wird \textbf{vorher} bestimmt, aus Basisrate, kleinstem
relevanten Effekt, Signifikanzniveau $\alpha$ und Power. Bei Basisrate 10\,\%:

\begin{center}
\begin{tabular}{lr}
\hline
\textbf{Zu erkennender Effekt} & \textbf{Nutzer je Gruppe} \\
\hline
$+20\,\%$ & 3\,841 \\
$+10\,\%$ & 14\,751 \\
$+5\,\%$  & 57\,763 \\
\hline
\end{tabular}
\end{center}

Halbiert man den Effekt, vervierfacht sich die Stichprobe.

\begin{notebox}
\textbf{Der p-Wert ist nicht die Wahrscheinlichkeit, dass die Änderung wirkt.} Er sagt:
wenn die Änderung nichts bewirkt, wie oft erzeugt der Zufall allein einen so großen
Unterschied.
\end{notebox}

\textbf{Peeking:} Wer den laufenden Test wiederholt anschaut und beim ersten
signifikanten Ergebnis abbricht, erhöht die Fehlerrate drastisch. In der Simulation
mit zwei \emph{identischen} Gruppen: 5\,\% falsche Signifikanzen bei einmaliger
Auswertung am Ende, 17\,\% bei zehnmaligem Hineinschauen.

% =========================================
\section{Architektur und Modellierung}
% =========================================

\subsection{Kopplung und Kohäsion}

\begin{definitionbox}
\textbf{Kopplung} ist das Maß dafür, wie stark ein Element mit anderen verbunden ist,
Wissen über sie besitzt oder von ihnen abhängt.

\textbf{Kohäsion} ist das Maß dafür, wie stark zusammenhängend die
Verantwortlichkeiten innerhalb eines Elements sind.
\end{definitionbox}

Die Entwurfsregel lautet: \textbf{Kopplung minimieren, Kohäsion maximieren}. Beides
zugleich ist oft ein Zielkonflikt, denn Verantwortlichkeiten zusammenzuziehen erhöht
die Kohäsion und schafft neue Beziehungen nach außen.

Typische Formen der Kopplung zwischen $A$ und $B$: $A$ hat ein Attribut vom Typ $B$,
$A$ ruft einen Dienst von $B$ auf, $A$ hat eine Methode mit $B$ als Parameter,
Rückgabewert oder lokaler Variable, $A$ erbt von $B$, $A$ implementiert die
Schnittstelle $B$.

Folgen starker Kopplung: lokale Änderungen ziehen Änderungen in verbundenen Klassen
nach sich, Klassen sind einzeln schwer zu verstehen und schwer wiederzuverwenden.
Folgen geringer Kohäsion: unklare Verantwortung, schlechte Verständlichkeit,
schlechtere Wartbarkeit.

\subsection{LCOM}
\label{sec:lcom}

\begin{definitionbox}
Für eine Klasse mit $a$ Feldern und $m$ Methoden sei $n(A_i)$ die Zahl der Methoden,
die auf Feld $A_i$ zugreifen. Dann ist
\[
LCOM = \frac{\left(\frac{1}{a}\sum_{i=1}^{a} n(A_i)\right) - m}{1 - m}
\]
mit Wertebereich $[0,1]$. $0$ bedeutet hohe Kohäsion (jede Methode nutzt jedes Feld),
$1$ bedeutet sehr geringe Kohäsion.
\end{definitionbox}

\subsubsection*{Beispiele}

\textbf{Dreieck:} $a=3$, $m=4$, jedes Feld von jeder Methode genutzt.
\[
LCOM = \frac{\frac{1}{3}\cdot 12 - 4}{1-4} = \frac{4-4}{-3} = 0
\]

\textbf{Figur (Dreieck und Kreis verschmolzen):} $a=4$, $m=8$, jedes Feld weiterhin von
nur vier Methoden genutzt.
\[
LCOM = \frac{\frac{1}{4}\cdot 16 - 8}{1-8} = \frac{4-8}{-7} = \frac{4}{7} \approx 0{,}57
\]

\textbf{Übungsaufgabe:} $A_1, A_2, A_3$ und $m_1 \dots m_4$ mit
$m_1: \{A_1, A_2\}$, $m_2: \{A_2\}$, $m_3: \{A_2, A_3\}$, $m_4: \{A_3\}$.
Also $n(A_1) = 1$, $n(A_2) = 3$, $n(A_3) = 2$, Summe $6$:
\[
LCOM = \frac{\frac{1}{3}\cdot 6 - 4}{1-4} = \frac{2-4}{-3} = \frac{2}{3} \approx 0{,}67
\]
Das ist weit von $0$ entfernt und deutet auf mittlere bis geringe Kohäsion hin.

\subsection{Sichten}

Eine Architektur wird nicht in einem Bild dargestellt, sondern in mehreren
\textbf{Sichten}, weil verschiedene Fragen verschiedene Strukturen betreffen:

\begin{center}
\begin{tabular}{lp{8.2cm}}
\hline
\textbf{Sicht} & \textbf{Beantwortet} \\
\hline
Baustein / Modul & Woraus besteht das System, wer hängt von wem ab? \\
Laufzeit         & Welche Prozesse und Objekte existieren zur Laufzeit, wie kommunizieren sie? \\
Verteilung       & Auf welcher Hardware läuft was? \\
\hline
\end{tabular}
\end{center}

\subsection{SOLID}

\begin{center}
\begin{tabular}{llp{7.2cm}}
\hline
\textbf{S} & Single Responsibility & Eine Klasse hat genau einen Grund, sich zu ändern \\
\textbf{O} & Open/Closed & Offen für Erweiterung, geschlossen für Änderung \\
\textbf{L} & Liskov Substitution & Untertypen müssen ohne Bruch an die Stelle des Obertyps treten können \\
\textbf{I} & Interface Segregation & Viele kleine Schnittstellen statt einer großen \\
\textbf{D} & Dependency Inversion & Abhängigkeit auf Abstraktionen, nicht auf Implementierungen \\
\hline
\end{tabular}
\end{center}

Single Responsibility ist die Kohäsionsregel aus Abschnitt~\ref{sec:lcom} in Worten,
Dependency Inversion die Kopplungsregel: Abhängigkeiten sollen von instabilen zu
stabilen Bausteinen zeigen.

\subsection{ADD 3.0: Attribute Driven Design}
\label{sec:add}

ADD ist ein \emph{iteratives} Entwurfsverfahren, das von Qualitätsattributen ausgeht
und nicht von Funktionen. Sieben Schritte je Iteration:

\begin{enumerate}
  \item \textbf{Eingaben prüfen} -- Ziele, Use Cases, Qualitätsszenarien,
        Randbedingungen und Concerns sichten und priorisieren.
  \item \textbf{Iterationsziel festlegen} -- welches Designziel diese Iteration
        verfolgt, etwa \glqq{}Primärfunktion strukturieren\grqq{} oder
        \glqq{}Latenzproblem adressieren\grqq{}.
  \item \textbf{Zu verfeinernde Elemente wählen} -- welche Systemteile bearbeitet
        werden, typischerweise vom Groben zum Detail.
  \item \textbf{Designkonzepte auswählen} -- Muster und Taktiken identifizieren und
        gegeneinander abwägen.
  \item \textbf{Elemente instanziieren, Verantwortungen zuordnen, Schnittstellen
        definieren} -- konkrete Bausteine samt Protokollen und Datenflüssen.
  \item \textbf{Sichten skizzieren und Entscheidungen protokollieren} -- inklusive
        Begründung und verworfener Alternativen.
  \item \textbf{Analyse und Review gegen das Iterationsziel} -- Abdeckung,
        Restrisiken, Aufgaben für die nächste Iteration.
\end{enumerate}

Iterativ heißt: nicht die ganze Architektur auf einmal, sondern je Runde ein
priorisiertes Ziel. Vorteile: frühe Rückmeldung, beherrschbare Entscheidungen und
die Möglichkeit, bei geänderten Treibern die Richtung zu wechseln.

% =========================================
\section{Automation}
% =========================================

\subsection{Build und Abhängigkeiten}

Ein Build-Werkzeug (Maven, Gradle, npm) macht drei Dinge reproduzierbar: die
\textbf{Abhängigkeiten} mit ihren Versionen, die \textbf{Schritte} vom Quelltext zum
Artefakt und die \textbf{Umgebung}, in der das geschieht.

Transitive Abhängigkeiten sind der schwierige Teil: A braucht B, B braucht C in
Version 1, D braucht C in Version 2. Version-Locking und Konfliktauflösung sind
deshalb Kernfunktionen des Werkzeugs und keine Nebensache.

\subsection{Continuous Integration und Delivery}

\begin{itemize}
  \item \textbf{Continuous Integration:} Jede Änderung wird zeitnah integriert und
        automatisch gebaut und getestet. Ziel ist, Integrationsfehler früh und in
        kleinen Portionen zu finden.
  \item \textbf{Continuous Delivery:} Jeder erfolgreiche Build ist prinzipiell
        auslieferbar. Die Freigabe bleibt eine Entscheidung.
  \item \textbf{Continuous Deployment:} Jeder erfolgreiche Build wird automatisch
        ausgeliefert.
\end{itemize}

Eine Pipeline ordnet die Prüfungen nach Kosten: erst Kompilieren und Modultests,
dann Integrationstests, dann langsame Prüfungen wie Last- oder Sicherheitstests. Wer
teuer vor billig prüft, bezahlt jede triviale Regression mit voller Laufzeit.

% =========================================
\section{Qualitätssicherung und Programmanalyse}
% =========================================

\subsection{Einordnung der Verfahren}

\begin{center}
\begin{tabular}{lll}
\hline
 & \textbf{statisch} (ohne Ausführung) & \textbf{dynamisch} (mit Ausführung) \\
\hline
\textbf{prüfend}  & Review, Coding Standards, Metriken & Test, Überdeckungsmessung \\
\textbf{beweisend}& Programmanalyse, Model Checking, Hoare & -- \\
\hline
\end{tabular}
\end{center}

Ein Test zeigt die Anwesenheit von Fehlern, nie deren Abwesenheit. Ein Beweis zeigt
die Abwesenheit, kostet aber eine Spezifikation und einen Aufwand, den man sich
gezielt aussuchen muss.

\subsection{Metriken}

\begin{itemize}
  \item \textbf{LOC:} Zeilen Code. Einfachster Indikator, misst Tippen und nicht
        Schwierigkeit.
  \item \textbf{LCOM:} fehlende Kohäsion, siehe Abschnitt~\ref{sec:lcom}.
  \item \textbf{CC:} zyklomatische Komplexität, siehe unten.
\end{itemize}

Ein \textbf{Coding Standard} ist erst dann ein Standard, wenn seine Regeln
maschinell prüfbar sind: ein Build kann daran scheitern. Werkzeuge wie PMD suchen
dafür Muster im abstrakten Syntaxbaum.

\subsection{Die Sprache While}

\[
S ::= [x := a]^l \mid [\texttt{skip}]^l \mid S_1; S_2
      \mid \texttt{if } [b]^l \texttt{ then } S_1 \texttt{ else } S_2
      \mid \texttt{while } [b]^l \texttt{ do } S
\]

Labels identifizieren Anweisungen \textbf{eindeutig}. Vier Funktionen darauf:

\[
\begin{aligned}
init([x := a]^l) &= l                & final([x := a]^l) &= \{l\} \\
init(S_1; S_2) &= init(S_1)          & final(S_1; S_2) &= final(S_2) \\
init(\texttt{if } [b]^l \dots) &= l  & final(\texttt{if } [b]^l \texttt{ then } S_1 \texttt{ else } S_2) &= final(S_1) \cup final(S_2) \\
init(\texttt{while } [b]^l \dots) &= l & final(\texttt{while } [b]^l \texttt{ do } S) &= \{l\}
\end{aligned}
\]

\[
\begin{aligned}
flow(S_1; S_2) &= flow(S_1) \cup flow(S_2) \cup \{(l, init(S_2)) : l \in final(S_1)\}\\
flow(\texttt{if } [b]^l \texttt{ then } S_1 \texttt{ else } S_2)
  &= flow(S_1) \cup flow(S_2) \cup \{(l, init(S_1)), (l, init(S_2))\}\\
flow(\texttt{while } [b]^l \texttt{ do } S)
  &= flow(S) \cup \{(l, init(S))\} \cup \{(l', l) : l' \in final(S)\}
\end{aligned}
\]

\begin{definitionbox}
Der \textbf{Kontrollflussgraph} eines Programms $S$ ist $(V, E)$ mit
$V = labels(S)$ und $E = flow(S)$.
\end{definitionbox}

\subsubsection*{Beispiel}

\begin{minipage}{0.42\textwidth}
\[
\begin{aligned}
&[y := x]^1;\\
&[z := 1]^2;\\
&\texttt{while } [y > 1]^3 \texttt{ do}\\
&\quad [z := x * y]^4;\\
&\quad [y := y - 1]^5;\\
&[y := 0]^6
\end{aligned}
\]
\end{minipage}
\hfill
\begin{minipage}{0.52\textwidth}
\begin{tikzpicture}[
  n/.style={draw, rounded corners, minimum width=1.9cm, minimum height=0.6cm,
            font=\scriptsize},
  t/.style={draw, diamond, aspect=2, inner sep=1pt, font=\scriptsize},
  >=Stealth, node distance=0.75cm
]
  \node[n] (a) {$[y := x]^1$};
  \node[n, below=of a] (b) {$[z := 1]^2$};
  \node[t, below=of b] (c) {$[y>1]^3$};
  \node[n, right=1.2cm of c] (d) {$[z := x{*}y]^4$};
  \node[n, below=of d] (e) {$[y := y{-}1]^5$};
  \node[n, below=of c] (f) {$[y := 0]^6$};

  \draw[->] (a) -- (b);
  \draw[->] (b) -- (c);
  \draw[->] (c) -- node[above, font=\tiny] {wahr} (d);
  \draw[->] (d) -- (e);
  \draw[->] (e) -- ++(1.3,0) |- ++(0,2.1) -| (c);
  \draw[->] (c) -- node[left, font=\tiny] {falsch} (f);
\end{tikzpicture}
\end{minipage}

\[
flow = \{(1,2), (2,3), (3,4), (4,5), (5,3), (3,6)\}
\]

\subsection{Zyklomatische Komplexität}

\begin{definitionbox}
$CC = e - n + 2p$ mit $e$ Kanten, $n$ Knoten und $p$ Zusammenhangskomponenten des
Kontrollflussgraphen. $CC$ ist die Zahl der unabhängigen Pfade.
\end{definitionbox}

Für das Beispiel: $6 - 6 + 2 \cdot 1 = 2$. Geradliniger Code hat $CC = 1$, jede
Verzweigung erhöht den Wert um eins. $CC$ ist eine untere Schranke für die Zahl der
Testfälle, die man für Zweigüberdeckung braucht.

\subsection{Reaching Definitions}

\begin{definitionbox}
Ein Paar $(x, l)$ bedeutet: die Zuweisung an Label $l$ kann hier noch der aktuelle
Wert von $x$ sein. $(x, ?)$ bedeutet: $x$ ist möglicherweise nicht initialisiert.
\end{definitionbox}

\[
\begin{aligned}
kill_{RD}([x := a]^l) &= \{(x,?)\} \cup \{(x,l') : [B]^{l'} \text{ weist } x \text{ zu}\} \\
gen_{RD}([x := a]^l) &= \{(x,l)\} \\
kill_{RD}([\texttt{skip}]^l) &= kill_{RD}([b]^l) = \emptyset,\quad
gen_{RD}([\texttt{skip}]^l) = gen_{RD}([b]^l) = \emptyset
\end{aligned}
\]

\[
RD_{entry}(l) =
\begin{cases}
\{(x,?) : x \in FV(S)\} & \text{für } l = init(S)\\
\bigcup \{RD_{exit}(l') : (l',l) \in flow(S)\} & \text{sonst}
\end{cases}
\]
\[
RD_{exit}(l) = \bigl(RD_{entry}(l) \setminus kill_{RD}([B]^l)\bigr) \cup gen_{RD}([B]^l)
\]

Die Analyse ist \textbf{vorwärts} und eine \textbf{may}-Analyse: vereinigt wird, weil
eine Definition ein Label erreicht, wenn sie es auf \emph{irgendeinem} Pfad erreicht.
Anwendung: Ist $(x,?)$ am Eingang eines Labels, das $x$ liest, liegt möglicherweise
eine Benutzung vor der Initialisierung vor.

\subsection{Die vier klassischen Datenflussanalysen}

\begin{center}
\begin{tabular}{lll}
\hline
 & \textbf{may} (Vereinigung) & \textbf{must} (Schnitt) \\
\hline
\textbf{vorwärts}  & Reaching Definitions & Available Expressions \\
\textbf{rückwärts} & Live Variables & Very Busy Expressions \\
\hline
\end{tabular}
\end{center}

Alle vier sind derselbe Algorithmus mit vier Stellschrauben: Richtung, Verknüpfung an
Zusammenflüssen, Startwert und die Funktionen $kill$ und $gen$. Die Transferfunktion
ist immer $(\text{eingehend} \setminus kill) \cup gen$.

\subsection{Strukturierte operationelle Semantik}

\begin{center}
\begin{tabular}{ll}
\hline
$\langle \texttt{skip}, v\rangle \Rightarrow v$ & [SKIP] \\
$\langle x := a, v\rangle \Rightarrow v[x \mapsto c]$, \ $c = \alpha[a,v]$ & [ASS] \\
$\langle \texttt{if } b \texttt{ then } S_1 \texttt{ else } S_2, v\rangle \Rightarrow \langle S_1, v\rangle$, falls $\beta[b,v]=1$ & [IF$_T$] \\
$\langle \texttt{if } b \texttt{ then } S_1 \texttt{ else } S_2, v\rangle \Rightarrow \langle S_2, v\rangle$, falls $\beta[b,v]=0$ & [IF$_F$] \\
$\langle \texttt{while } b \texttt{ do } S, v\rangle \Rightarrow \langle S; \texttt{while } b \texttt{ do } S, v\rangle$, falls $\beta[b,v]=1$ & [WH$_T$] \\
$\langle \texttt{while } b \texttt{ do } S, v\rangle \Rightarrow v$, falls $\beta[b,v]=0$ & [WH$_F$] \\
$\dfrac{\langle S_1, v\rangle \Rightarrow \langle S_1', v'\rangle}{\langle S_1;S_2, v\rangle \Rightarrow \langle S_1';S_2, v'\rangle}$ & [COMP$_1$] \\[10pt]
$\dfrac{\langle S_1, v\rangle \Rightarrow v'}{\langle S_1;S_2, v\rangle \Rightarrow \langle S_2, v'\rangle}$ & [COMP$_2$] \\
\hline
\end{tabular}
\end{center}

Zwei Axiome überführen eine Konfiguration in einen \emph{Zustand} (SKIP, ASS, WH$_F$),
die übrigen in eine \emph{Konfiguration}. Die Schleife wird pro Schritt einmal
abgerollt; deshalb beschreibt ein endliches Regelsystem eine unbeschränkte Ausführung.

\subsection{Symbolische Ausführung}

Dieselben Regeln, aber die Variablen zeigen auf \emph{Ausdrücke über symbolischen
Werten}, und jeder Zweig trägt eine \textbf{Pfadbedingung} $\varphi$:

\[
\langle \texttt{if } b \texttt{ then } S_1 \texttt{ else } S_2, v, \varphi\rangle
\Rightarrow \langle S_1, v, \varphi \wedge \sigma[b,v]\rangle
\quad\text{falls erfüllbar}
\]

Beide Regeln können gelten, es entsteht ein \textbf{Ausführungsbaum}. Ein Zweig mit
unerfüllbarer Pfadbedingung wird verworfen; das Lösen einer Pfadbedingung liefert
konkrete Eingaben, die genau dieses Blatt erreichen -- also automatisch erzeugte
Testfälle.

Grenze: jede Schleifeniteration ist eine weitere Abrollung. Bei unbeschränkter
Iterationszahl ist der Baum unendlich, und eine Tiefenschranke ist die praktische
Antwort, keine Lösung.

% =========================================
\section{Verifikation}
% =========================================

\subsection{DPLL}

Entscheidet die Erfüllbarkeit einer Formel in KNF. Zwei Beobachtungen sparen ganze
Teilbäume:

\begin{itemize}
  \item \textbf{Unit-Klausel:} Bleibt in einer Klausel genau ein unbelegtes Literal,
        ist seine Belegung erzwungen.
  \item \textbf{Pures Literal:} Kommt eine Variable nur mit einem Vorzeichen vor, kann
        sie entsprechend belegt werden, ohne je zu schaden.
\end{itemize}

\begin{center}
\begin{minipage}{0.92\textwidth}
\small
\texttt{DPLL(C, v):}\\
\hspace*{1em}$C^- \leftarrow \{c \in C : c[v] \to false\}$;
\textbf{if} $C^-$ nicht leer \textbf{then return} $false$\\
\hspace*{1em}$C^\star \leftarrow C \setminus \{c \in C : true \to c[v]\}$;
\textbf{if} $C^\star$ leer \textbf{then return} $true$\\
\hspace*{1em}\textbf{if} Unit-Klausel für $x$ oder $\bar{x}$ \textbf{then return}
$DPLL(C^\star, v[x \mapsto b])$\\
\hspace*{1em}\textbf{if} pures Literal $x$ oder $\bar{x}$ \textbf{then return}
$DPLL(C^\star, v[x \mapsto b])$\\
\hspace*{1em}wähle $x$; \textbf{return} $DPLL(C^\star, v[x \mapsto 0]) \vee DPLL(C^\star, v[x \mapsto 1])$
\end{minipage}
\end{center}

\subsubsection*{Beispiel}

$(x_1 \vee x_2) \wedge (\neg x_2 \vee x_3) \wedge (x_2 \vee \neg x_3)$

\begin{enumerate}
  \item $x_1$ ist pur $\Rightarrow v[x_1 \mapsto 1]$, erste Klausel erfüllt
  \item Verzweigung $v[x_2 \mapsto 0]$
  \item $\neg x_3$ ist Unit $\Rightarrow v[x_3 \mapsto 0]$, $C^\star$ leer
\end{enumerate}

Eine geprüfte Belegung statt $2^3 = 8$. Der Worst Case bleibt $O(2^n)$: SAT ist
NP-vollständig, und keine der beiden Regeln ändert daran etwas.

\subsection{Transitionssysteme}

\begin{definitionbox}
Ein \textbf{Transitionssystem} ist ein Tripel $T = \langle s, I, \Gamma\rangle$ mit
einer Menge getypter Variablen $s$, einem Booleschen Ausdruck $I$ über $s$ für die
Anfangszustände und einem Ausdruck $\Gamma$ über $s \cup s'$ für die
Transitionsrelation.
\end{definitionbox}

Eine Komponente hat Eingaben und Ausgaben. Komponiert man sie mit einem
\textbf{Treiber}, der Eingaben liefert, und einem \textbf{Monitor}, der Ausgaben
prüft, entsteht ein geschlossenes System ohne Ein- und Ausgaben -- genau ein
Transitionssystem. Komposition ist Konjunktion der Relationen; gemeinsame Namen sind
die Verdrahtung.

\subsection{Bounded Model Checking}

\begin{definitionbox}
Mit $\Gamma_{i,j} = \Gamma[s_i/s][s_j/s']$ und $P^k = P[s_k/s]$ heißt $T$
\textbf{$k$-sicher} für $P$, wenn
\[
\bigl(I[s_0/s] \wedge \Gamma_{0,1} \wedge \dots \wedge \Gamma_{k-1,k}\bigr)
\models \bigl(P^0 \wedge \dots \wedge P^k\bigr)
\]
\end{definitionbox}

Getestet wird wie jede Implikation, über ein Modell der Negation:
\[
I[s_0/s] \wedge \Gamma_{0,1} \wedge \dots \wedge \Gamma_{k-1,k}
\wedge \bigl(\neg P^0 \vee \dots \vee \neg P^k\bigr)
\]

\textbf{Algorithmus:} $k = 0$; ist der Test erfüllbar, ist ein Fehler gefunden und das
Modell ist das Gegenbeispiel; sonst $k \leftarrow k+1$ und erneut.

\subsubsection*{Beispiel (Klausuraufgabe)}

$T = \langle \{s_1\},\ s_1 = 20,\ s_1' = s_1 / 2 \rangle$, $P = (s_1 \bmod 2 = 0)$.

\[
\Gamma_{0,1}: s_1^{(1)} = s_1^{(0)}/2, \quad
\Gamma_{1,2}: s_1^{(2)} = s_1^{(1)}/2, \quad
\Gamma_{2,3}: s_1^{(3)} = s_1^{(2)}/2
\]

Der Lauf ist $20 \to 10 \to 5 \to 2$. Für $k=0$ und $k=1$ ist das System $k$-sicher;
für $k=2$ ist der Test erfüllbar, das Gegenbeispiel ist $20, 10, 5$ und verletzt $P$
im dritten Zustand.

\begin{notebox}
BMC beweist \textbf{nie} Sicherheit. Die Menge der in $i$ Schritten erreichbaren
Zustände wächst mit $i$, und es gibt keine Formel, die sagt, wann sie aufhört zu
wachsen. Genau diese Lücke schließen induktive Invarianten.
\end{notebox}

\subsection{Induktive Invarianten}

\begin{definitionbox}
$P$ ist eine \textbf{induktive Invariante} von $T = \langle s, I, \Gamma\rangle$, wenn
\begin{enumerate}
  \item \textbf{Basis:} $I \models P$
  \item \textbf{Schritt:} $P \wedge \Gamma \models P[s'/s]$
\end{enumerate}
\end{definitionbox}

Gelten beide, so gilt $P$ in jedem erreichbaren Zustand, per Induktion über die
Länge des Laufs -- ohne jede Schranke.

\textbf{Der übliche Stolperstein:} Der Schrittbeweis startet in \emph{jedem} Zustand,
der $P$ erfüllt, auch in unerreichbaren. Beispiel: $x' = x + 2$ mit $I: x = 0$ und
$P: x \neq 5$. $P$ gilt in allen erreichbaren Zuständen (sie sind gerade), ist aber
nicht induktiv: aus $x = 3$ folgt $x' = 5$. $x = 3$ ist unerreichbar, und der
Schritttest weiß das nicht.

\textbf{Verstärkung:} Man sucht ein stärkeres $Q$, das $P$ impliziert und induktiv
ist. Hier genügt $x \bmod 2 = 0$, und $P \wedge (x \bmod 2 = 0)$ ist induktiv.

\subsection{Hoare-Logik}

\begin{definitionbox}
Das Tripel $\{P\}\, S\, \{Q\}$ bedeutet: gilt $P$ vor $S$ und terminiert $S$, so gilt
$Q$ danach (\emph{partielle} Korrektheit).
\end{definitionbox}

\begin{center}
\begin{tabular}{ll}
\hline
$\{Q\}\ \texttt{skip}\ \{Q\}$ & Skip \\
$\{Q[a/x]\}\ x := a\ \{Q\}$ & Zuweisung \\
$\dfrac{\{P\} S_1 \{R\},\ \{R\} S_2 \{Q\}}{\{P\}\ S_1;S_2\ \{Q\}}$ & Komposition \\[10pt]
$\dfrac{\{P \wedge b\} S_1 \{Q\},\ \{P \wedge \neg b\} S_2 \{Q\}}{\{P\}\ \texttt{if } b \texttt{ then } S_1 \texttt{ else } S_2\ \{Q\}}$ & Verzweigung \\[10pt]
$\dfrac{\{I \wedge b\}\ S\ \{I\}}{\{I\}\ \texttt{while } b \texttt{ do } S\ \{I \wedge \neg b\}}$ & Schleife \\[10pt]
$\dfrac{P \to P',\ \{P'\} S \{Q'\},\ Q' \to Q}{\{P\}\ S\ \{Q\}}$ & Konsequenz \\
\hline
\end{tabular}
\end{center}

Die Zuweisungsregel läuft \textbf{rückwärts}: um zu wissen, was vor $x := a$ gelten
muss, ersetzt man $x$ in $Q$ durch $a$. Vorwärts wäre es geraten, rückwärts ist es
eine Gleichung.

Die Schleife ist die einzige Regel mit echter Arbeit: sie braucht eine
\textbf{Invariante}, die keine Regel liefert.

\subsubsection*{Beispiel (Übungsblatt 6)}

\begin{center}
\begin{minipage}{0.62\textwidth}
\small
\texttt{\{ true \}}\\
\texttt{if (a >= b) \{ while (a != b) a := a - 1; \}}\\
\texttt{else \{ while (a != b) a := a + 1; \}}\\
\texttt{\{ a = b \}}
\end{minipage}
\end{center}

Invarianten: im \texttt{then}-Zweig $a \geq b$, im \texttt{else}-Zweig $a \leq b$.
Nachweis für den \texttt{then}-Zweig: aus $a \geq b \wedge a \neq b$ folgt $a > b$,
also $a - 1 \geq b$, die Invariante bleibt erhalten. Beim Verlassen gilt
$a \geq b \wedge \neg(a \neq b)$, also $a = b$.

% =========================================
\section{Modellbasiertes Testen}
% =========================================

\subsection{Mealy-Automaten}

Ein Mealy-Automat antwortet auf jedes Eingabesymbol mit einem Ausgabesymbol; ein Wort
der Länge $n$ erzeugt ein Ausgabewort der Länge $n$. Damit ist er das natürliche
Modell einer reaktiven Komponente.

\begin{center}
\begin{tabular}{ccc}
\hline
\textbf{Zustand} & \textbf{bei a} & \textbf{bei b} \\
\hline
$q_0$ & $q_1/0$ & $q_2/1$ \\
$q_1$ & $q_1/1$ & $q_2/0$ \\
$q_2$ & $q_0/0$ & $q_2/1$ \\
\hline
\end{tabular}
\end{center}

Für $M_1$: $Out(abba) = 0010$ und $Out(baab) = 1000$. Ein \textbf{Testorakel}
vergleicht die behauptete Ausgabe eines Systems mit dieser Vorhersage.

\subsection{Überdeckungen}

\begin{center}
\begin{tabular}{lll}
\hline
\textbf{Überdeckung} & \textbf{Idee} & \textbf{Größe für $M_1$} \\
\hline
Zustandsüberdeckung & jeden Zustand erreichen & 3 \\
Transitionsüberdeckung & jede Transition ausführen & 6 \\
Transition-Switch & jedes ausführbare Transitionspaar & 12 \\
W-Methode & Transitionsüberdeckung + Unterscheidungsmenge & vollständig unter Annahme \\
\hline
\end{tabular}
\end{center}

Bei $n$ Zuständen und $k$ Eingabesymbolen gibt es $n \cdot k \cdot k$ Paare direkt
aufeinanderfolgender Transitionen, für $M_1$ also $3 \cdot 2 \cdot 2 = 12$. Die Suite
baut man systematisch über Zugriffssequenzen $\alpha(q)$:
\[
T = \{\alpha(q)\,xy : q \in Q,\ x,y \in \Sigma\}
\]
mit $\alpha(q_0) = \varepsilon$, $\alpha(q_1) = a$, $\alpha(q_2) = b$.

\subsection{Observation Table und L*}

Ein Lerner sieht das System nicht und darf zwei Fragen stellen:
\textbf{Membership Query} (führe dieses Wort aus, was kommt zurück?) und
\textbf{Equivalence Query} (stimmt diese Hypothese, und wenn nicht, wo weicht sie ab?).

\begin{definitionbox}
Die \textbf{Observation Table} hat Zeilen $U$ (Zugriffswörter) und $U\!\cdot\!A$
(Rand), Spalten $V$ (Suffixe) und in jeder Zelle das \emph{letzte} Ausgabesymbol von
$u \cdot v$.
\begin{itemize}
  \item \textbf{closed:} jede Zeile aus $U\!\cdot\!A$ kommt bereits in $U$ vor
  \item \textbf{consistent:} gleich aussehende Zeilen bleiben nach einem weiteren
        Symbol gleich
\end{itemize}
\end{definitionbox}

Reparatur: nicht geschlossen $\Rightarrow$ Zeile nach $U$ verschieben (neuer Zustand);
nicht konsistent $\Rightarrow$ Spalte hinzufügen (neues Experiment).

\textbf{Ablauf von L*:} Tabelle füllen, schließen, konsistent machen, Hypothese bauen,
Äquivalenz anfragen; bei Gegenbeispiel Tabelle verfeinern und wiederholen. Jede Runde
fügt mindestens einen Zustand hinzu, das Ziel hat endlich viele, also terminiert es.

\begin{notebox}
Die Equivalence Query beantwortet in der Praxis niemand. Sie wird durch eine
Testsuite angenähert -- und damit ist das gelernte Modell nur so gut wie diese Suite.
Lernen liefert die Tests, die Tests liefern das Lernen: dieser Kreis ist der
tatsächliche Stand des Verfahrens, kein Darstellungsfehler.
\end{notebox}

% =========================================
\section{Formelsammlung}
% =========================================

\begin{center}
\begin{tabular}{lp{9.6cm}}
\hline
\textbf{Thema} & \textbf{Formel} \\
\hline
Zyklomatische Komplexität & $CC = e - n + 2p$ \\
LCOM & $\left(\frac{1}{a}\sum n(A_i) - m\right) / (1-m)$, Bereich $[0,1]$ \\
Punktschätzung & $E = \frac{a + r c + b}{2+r}$, \ $S = \frac{b-a}{u}$;
                 2-Pkt.: $r=0$, 3-Pkt./PERT: $r=4$, $u=6$ \\
Aggregation & $E = \sum E_i$, \ $S = \sqrt{\sum S_i^2}$ \\
Velocity & Story Points fertiger Stories pro Sprint; Planung mit Median \\
Implikation & $\varphi \models \psi$ gdw. $\varphi \wedge \neg\psi$ unerfüllbar \\
$k$-Sicherheit & $I[s_0] \wedge \Gamma_{0,1} \wedge \dots \wedge \Gamma_{k-1,k}
                  \models P^0 \wedge \dots \wedge P^k$ \\
Induktive Invariante & Basis $I \models P$; Schritt $P \wedge \Gamma \models P[s'/s]$ \\
Hoare-Zuweisung & $\{Q[a/x]\}\ x := a\ \{Q\}$ \\
Reaching Definitions & $RD_{exit}(l) = (RD_{entry}(l) \setminus kill) \cup gen$ \\
Transitionspaare & $n \cdot k^2$ bei $n$ Zuständen und $k$ Eingabesymbolen \\
\hline
\end{tabular}
\end{center}

\vspace{1em}
\begin{notebox}
Zu jedem Kapitel dieser Notizen gibt es lauffähigen Code im Verzeichnis
\texttt{swk/}: \texttt{program-analysis} für While, CFG und Datenfluss,
\texttt{verification} für DPLL, BMC, Invarianten und Hoare,
\texttt{model-based-testing} für Mealy, L* und Überdeckungen,
\texttt{metrics} für LCOM und Kopplung, \texttt{estimation} für Schätzung, Velocity
und A/B-Tests. Alle Beispiele dieser Notizen sind dort nachgerechnet.
\end{notebox}

\end{document}
